\documentclass{article}
\linespread{1.3}
\usepackage[margin=50pt]{geometry}
\usepackage{amsmath, amsthm, amssymb, amsthm, tikz, fancyhdr, float, graphicx, spverbatim, systeme}
\pagestyle{fancy}
\renewcommand{\headrulewidth}{0pt}
\newcommand{\changefont}{\fontsize{15}{15}\selectfont}

\fancypagestyle{firstpageheader}{
  \fancyhead[R]{
    \changefont
    \parbox[t]{4cm}{ % Adjust width as needed
      Michael Huang\\
      EN.555.644.81\\
      Homework 10
    }
  }
}

\begin{document}

\thispagestyle{firstpageheader}
{\Large 

\section*{14.2}
A company's cash position, measured in millions of dollars, follows a generalized Wiener process with a drift rate of 0.5 per quarter and a variance rate of 4.0 per quarter. How high does the company's initial cash position have to be for the company to have a less than 5\% chance of a negative cash position by the end of one year? \\ \\
We have drift of 0.5 and variance of 4.0 per quarter, and we calculate over a year, so 4 quarters. The cash position therefore has a normal distribution with \(\Phi(x + a \Delta t, b^2 \Delta t) = \Phi(x + 0.5 \cdot 4, 4 \cdot 4) = \Phi(x + 2, 16)\), so to get the "less than 5\% chance of negative cash position", we simply need to calculate 
\[
  N(-\frac{x + 2}{\sqrt{16}}) = 0.05
\]
\[
  N(-\frac{x + 2}{4}) = 0.05
\]
Looking up in a z score table, taking the approximation:
\[
  -\frac{x + 2}{4} = -1.64
\]
\[
  x + 2 = 6.56
\]
\[
  x = 4.56
\]
So the company's initial cash position would have to be approximately \boxed{\textbf{\$4.56 million}}


}
{\Large 

\section*{14.4}
Consider a variable, \(S\), that follows the process 
\[
  dS = \mu dt + \sigma dz
\]
For the first three years, \(\mu = 2\) and \(\sigma = 3\); for the next three years, \(\mu = 3\) and \(\sigma = 4\). If the initial value of the variable is 5, what is the probability distribution of the value of the variable at the end of year six? \\ \\
For the first three years, we have a normal distribution of \(\Phi(2 \cdot 3, 3^2 \cdot 3) = \Phi(6, 27)\); for the next three years, we have a normal distribution of \(\Phi(3 \cdot 3, 4^2 \cdot 3) = \Phi(9, 48)\). Due to the nature of the normal distribution, with the non-overlapping time intervals, we can find the total change by simply taking the sum of the independent normal random variables, so the total change is therefore \(\Phi(6 + 9, 27 + 48) = \Phi(15, 75)\). With the initial value, we therefore have a final probability distribution of \(\Phi(5 + 15, 75) = \boxed{\mathbf{\Phi(20, 75)}}\).

}
{\Large 

\section*{14.8}
It has been suggested that the short-term interest rate, \(r\), follows the stochastic process
\[
  dr = a(b - r)dt + rc dz
\]
where \(a, b\), and \(c\) are positive constants and \(dz\) is a Wiener process. Describe the nature of this process. \\ \\
From looking at the equation, we can see that the drift rate \(a\) is \(a(b-r)\) and the volatility \(b\) is \(rc\). From looking just at the drift rate, due to the \((b-r)\) component, when \(b < r\), or that the interest rate is above \(b\), we have negative drift rate, while we have positive drift rate when the opposite is true. Considering this negative/positive switchover point at \(b\), we see that the drift rate will continuously be pulled to be centered on \(b\) with a factor of \(a\). Looking at the volatility term, we see that there is a constant volatility of \(c\) applied to the change in the rate, with it then being multiplied by \(r\). Overall, we also see that the added change  in the rate due to volatility also increases as the rate is higher, and is lower when the it is lower. How severe this change is depends on the constant \(c\).

}
{\Large 

\section*{14.10}
Suppose that \(x\) is the yield to maturity with continuous compounding on a zero-coupon bond that pays off \$1 at time \(T\). Assume that \(x\) follows the process
\[
  dx = a(x_0 - x)dt + sx dz
\]
where \(a, x_0\), and \(s\) are positive constants and \(dz\) is a Wiener process. What is the process followed by the bond price? \\ \\
To determine the process for calculating the bond price \(P\), we first use the formula for the price of a zero-coupon bond with continuous discounting, which gets us \(P(x, t) = e^{-x(T - t)}\). As \(P\) is a function of the stochastic variable \(x\) and time \(t\), we can use Ito's Lemma here. We start by computing the corresponding partial derivatives: \\
- \(\frac{\partial P}{\partial x} = e^{-x(T-t)} \cdot -(T-t) = -(T-t)P(x, t)\) \\
- \(\frac{\partial P}{\partial t} = e^{-x(T-t)} \cdot x = x P(x, t)\) \\
- \(\frac{\partial^2 P}{\partial x^2} = e^{-x(T-t)} \cdot -(T-t) \cdot -(T-t) = (T-t)^2 e^{-x(T-t)} = (T - t)^2 P(x,t)\) \\
which we plug into the equation from Ito's Lemma, applied to our \(dx\):
\[
  dP = [\frac{\partial P}{\partial x} a + \frac{\partial P}{\partial t} + \frac{1}{2} \frac{\partial^2 P}{\partial x^2} b^2]dt + \frac{\partial P}{\partial x} b dz
\]
\[
  = [\frac{\partial P}{\partial x} a(x_0 - x) + \frac{\partial P}{\partial t} + \frac{1}{2} \frac{\partial^2 P}{\partial x^2} (sx)^2]dt + \frac{\partial P}{\partial x} dx dz
\]
\[
  = [-(T-t)P \cdot a(x_0 - x) + xP + \frac{1}{2} (T-t)^2 \cdot P (sx)^2]dt + -(T - t) \cdot P sx dz
\]
\[
  \boxed{\mathbf{dP = [-a (x_0 - x) (T-t) + x + \frac{1}{2} s^2x^2 (T-t)^2]Pdt + -sx(T - t) P dz}}
\]
We could also factor out P to further simplify the evaluation of the function. 

}
{\Large 

\section*{14.12}
Suppose that a stock price has an expected return of 16\% per annum and a volatility of 30\% per annum. When the stock price at the end of a certain day is \$50, calculate the following:

\subsection*{(a)} 
The expected stock price at the end of the next day. \\ \\
We know that 
\[
  \frac{dS}{S} = \mu dt + \sigma dz
\]
and we have \(\mu = 0.16, \sigma = 0.30\). We also have \(S = 50\) and \(\Delta t = \frac{1}{252}\) as we only consider trading days. Due to the nature of the Wiener process we have 
\[
  \frac{dS}{S} = \mu dt + \sigma dz  = \Phi(\mu \Delta t, \sigma^2 \Delta t)
\]
\[
  \frac{dS}{S} = \Phi(0.16 \cdot \frac{1}{252}, 0.30^2 \cdot \frac{1}{252})
\]
\[
  \frac{dS}{50} = \Phi(0.00063492063, 0.00035714285)
\]
\[
  dS = \Phi(0.00063492063 \cdot 50, 0.00035714285 \cdot 50^2)
\]
\[
  dS = \Phi(0.03174603174, 0.892857125)
\]
So \(dS\) for the next day has \(\mu = 0.03174603174\), so the next day's expected price is \(50 + 0.03174603174 = 50.03174603174 \approx\) \boxed{\textbf{\$50.03}}.

\subsection*{(b)}
The standard deviation of the stock price at the end of the next day. \\ \\
According to the distribution we found in the previous part, the standard deviation is \(\sqrt{0.892857125} = \boxed{\mathbf{0.94491117307}}\).

\subsection*{(c)}
The 95\% confidence limits for the stock price at the end of the next day. \\ \\
We take the 95\% CI according to the distribution:
\[
  [50.03 - 1.96 \cdot 0.94491117307, 50.03 + 1.96 \cdot 0.94491117307]
\]
\[
  [50.02 - 1.85202589922, 50.02 + 1.85202589922]
\]
\[
  [48.1679741008, 51.8720258992]
\]
So the 95\% confidence limits for the stock price are approximately \boxed{\textbf{[\$48.17, \$51.87]}}.

}

\end{document}